\Askhsh{22407}{4}
\begin{ylh}{1}
    \item 7.6. Διαίρεση τμημάτων εσωτερικά και εξωτερικά ως προς δοσμένο λόγο
    \item 7.7. Θεώρημα του Θαλή
    \item 9.2. Το Πυθαγόρειο θεώρημα
    \item 10.3. Εμβαδόν βασικών ευθύγραμμων σχημάτων
    \item 10.5. Λόγος εμβαδών όμοιων τριγώνων - πολυγώνων.
\end{ylh}
\begin{wrapfigure}[11]{r}{6cm}
\centering
\begin{tikzpicture}[scale=0.5] % Adjusted scale for better fit
    \tkzDefPoint(0,0){K}
    \tkzDefPoint(9,0){A}
    \tkzDefPoint(19,0){B} % Since AB=10, A is at 9, B is at 19
    \tkzDefPoint(0,12){G} % Gamma, from Pythagoras GK=12
    \tkzDefPoint(11.4,4.8){D} % Delta, calculated from Angle Bisector Theorem for D on BC
    \tkzDefPoint(11.4,0){L} % Lambda, projection of D

    % Draw the main line segment KB
    \tkzDrawSegment[line width=0.8pt](K,B) % Draw the line K-A-L-B as a continuous segment

    % Draw segments of the triangle and bisector/projections
    \tkzDrawSegments[line width=0.8pt](G,A G,B A,D D,B D,L G,K)

    % Mark right angles
    \tkzMarkRightAngle[fill=maincolor!20,draw=maincolor](G,K,A)
    \tkzMarkRightAngle[fill=maincolor!20,draw=maincolor](D,L,A)

    % Mark angle bisector
    \tkzMarkAngle[arc=ll,size=0.6cm,fill=maincolor!20,draw=maincolor](G,A,D)
    \tkzMarkAngle[arc=ll,size=0.6cm,fill=maincolor!20,draw=maincolor](D,A,B)

    % Draw points
    \tkzDrawPoints[color=black, fill=black](K,A,L,B,G,D)

    % Labels
    \tkzLabelPoint[below](K){K}
    \tkzLabelPoint[below](A){A}
    \tkzLabelPoint[below](L){$\Lambda$}
    \tkzLabelPoint[below](B){B}
    \tkzLabelPoint[above left](G){$\Gamma$}
    \tkzLabelPoint[above right](D){$\Delta$}
\end{tikzpicture}
\end{wrapfigure}
ΘΕΜΑ 4

Στο παρακάτω σχήμα η $A\Delta$ είναι διχοτόμος του τριγώνου $AB\Gamma$ και η $AK$ είναι η προβολή της πλευράς $A\Gamma$ πάνω στην ευθεία $AB$. Δίνονται $AB = 10$, $A\Gamma = 15$ και $AK = 9$.
\begin{alist}
    \item Να αποδείξετε ότι:
    \begin{rlist}
        \item $\Gamma K = 12$ και $(\text{AB}\Gamma) = 60$. \monades{8}
        \item $(\text{A}\Delta\text{B}) = 24$ και $(\text{A}\Delta\Gamma) = 36$. \monades{10}
    \end{rlist}
    \item Έστω $\Lambda$ η προβολή του σημείου $\Delta$ πάνω στην ευθεία $AB$.
    \begin{rlist}
        \item Να αποδείξετε ότι $\displaystyle \frac{\Delta\Lambda}{\Gamma K} = \frac{2}{5}$. \monades{3}
        \item Να βρείτε τον λόγο $\displaystyle \frac{\Lambda B}{\Lambda K}$ στον οποίο το σημείο $\Lambda$ διαιρεί εσωτερικά το ευθύγραμμο τμήμα $BK$. \monades{4}
    \end{rlist}
\end{alist}